#######################################################
PROMPT
#######################################################

all the calculations for the API for composition -> conductivity value should be in conductivity/
  projected_analytical_conductivity.py , use electrolyte_model.py , electrochem_tools.py , conductivity/electrolyte_utils_features.py ,
  constants.py , utils/ , data/species_data.py and species_fns.py

  code the full model in conductivity/projected_analytical_conductivity.py, and make the
  tests for the baseline electrolyte EC:DMC 3:7 v/v + 1M LiPF6 should assert within 8-10 mS/cm, electrolyte_property_db.py formulations (report 
  error but model is in the right direction if the conductivity predictions are +/- 1 mS/cm), 
  and check for all the physical tests i'll provide below :

  \documentclass[11pt]{article}

  \usepackage{amsmath,amssymb,mathtools,bm,booktabs,longtable}
  \usepackage[margin=1in]{geometry}

  \newcommand{\dd}{\,\mathrm d}
  \newcommand{\R}{\mathbb R}
  \newcommand{\E}{\mathbb E}
  \newcommand{\Tr}{\operatorname{Tr}}
  \newcommand{\diag}{\operatorname{diag}}
  \newcommand{\rank}{\operatorname{rank}}
  \newcommand{\argmax}{\operatorname*{arg\,max}}
  \newcommand{\argmin}{\operatorname*{arg\,min}}
  \newcommand{\one}{\mathbf 1}
  \newcommand{\zero}{\mathbf 0}
  \newcommand{\eps}{\varepsilon}

  \begin{document}

  \section*{Complete Projected Analytical Conductivity Model}

  \subsection*{Purpose}

  The model maps one electrolyte composition \(x\) to one scalar conductivity
  \[
  \sigma(x)\quad [\mathrm{S\,m^{-1}}],
  \]
  or equivalently
  \[
  \sigma_{\mathrm{mS/cm}}(x)=10\,\sigma(x).
  \]

  No Green--Kubo time integral is evaluated in production.
  No Einstein--Helfand long-time slope is evaluated in production.
  The model computes conductivity from projected equilibrium, flux, displacement,
  self-current, and memory primitives.

  \[
  x
  \longmapsto
  (U_x,D_x,P_x,\mathcal A,\Psi)
  \longmapsto
  (c,K,Q,d,M,D^{\rm self},A,h)
  \longmapsto
  \sigma.
  \]

  \section{Notation and operation glossary}

  \subsection{Vectors, matrices, tensors}

  A vector is a column vector unless stated otherwise.

  For vectors \(u,v\in\R^n\),
  \[
  u^\top v=\sum_{i=1}^n u_i v_i.
  \]

  For a square matrix \(B\),
  \[
  \Tr(B)=\sum_i B_{ii}.
  \]

  For a symmetric matrix \(B\),
  \[
  B\succeq0
  \]
  means
  \[
  y^\top B y\ge0
  \qquad
  \text{for all }y.
  \]

  For a tensor \(T_i\in\R^{3\times3}\),
  \[
  \Tr(T_i)=T_{i,xx}+T_{i,yy}+T_{i,zz}.
  \]

  \subsection{The symbol \(B^\#\)}

  The symbol
  \[
  B^\#
  \]
  means the Moore--Penrose pseudoinverse when \(B\) is a symmetric positive
  semidefinite matrix.

  Implementation:

  1. Compute the eigendecomposition
  \[
  B=V\Lambda V^\top,
  \]
  where
  \[
  \Lambda=\diag(\lambda_1,\ldots,\lambda_m),
  \qquad
  \lambda_r\ge0.
  \]

  2. Choose numerical tolerance
  \[
  \tau_{\#}=10^{-12}\max_r \lambda_r.
  \]

  3. Define
  \[
  \lambda_r^\#=
  \begin{cases}
  1/\lambda_r, & \lambda_r>\tau_{\#},\\
  0, & \lambda_r\le\tau_{\#}.
  \end{cases}
  \]

  4. Return
  \[
  B^\#=V\diag(\lambda_1^\#,\ldots,\lambda_m^\#)V^\top.
  \]

  If a vector \(y\) has nonzero projection into the nullspace of \(B\), then
  \[
  y^\top B^\# y
  \]
  silently discards that null component. For a physical memory matrix this is
  allowed only if the discarded component is numerically zero. Therefore require
  \[
  \|(I-BB^\#)y\| \le \tau_{\rm null}\|y\|
  \]
  for each current-coupling vector \(y\), with for example
  \[
  \tau_{\rm null}=10^{-8}.
  \]

  For the finite Markov generator \(Q\), do not directly form \((-Q)^\#\).
  Instead solve the constrained Poisson system defined in Section 11.

  \section{Composition input}

  The electrolyte composition is
  \[
  x=
  \left(
  T,\,
  \{C_a\}_{a\in\mathcal C},\,
  \{\phi_s\}_{s\in\mathcal S},\,
  \mathcal P
  \right).
  \]

  Definitions:

  \[
  T>0
  \]
  is temperature in K.

  \[
  \mathcal C
  \]
  is the set of conserved components, for example Li, PF6, FSI, solvent molecules,
  neutral additives, or coarse-grained site classes.

  \[
  C_a
  \quad [\mathrm{mol\,m^{-3}}]
  \]
  is the analytical total concentration of conserved component \(a\).

  \[
  \phi_s
  \]
  is the liquid fraction of solvent/additive species \(s\).

  \[
  \mathcal P
  \]
  is the physical parameter map needed to construct the generator:
  \[
  \mathcal P
  =
  \left(
  \text{charge sites},
  \text{topology},
  \text{interaction parameters},
  \text{mobility parameters},
  \text{basin definitions},
  \text{memory coordinate definitions}
  \right).
  \]

  A recipe containing only names such as EC, DMC, LiPF6, FEC, VC, FSI, and
  concentrations is not mathematically enough. The model requires enough
  information to build the following functions:

  \[
  U_x(q),
  \qquad
  D_x(q),
  \qquad
  P_x(q),
  \qquad
  A_i\subset\Omega_x,
  \qquad
  \psi_\alpha(q).
  \]

  \section{Configuration space}

  Let
  \[
  q\in\Omega_x\subseteq\R^d
  \]
  be the model coordinate.

  For atomistic coordinates:
  \[
  q=(r_1,\ldots,r_N),
  \qquad
  r_a\in\R^3.
  \]

  For a reduced analytical model:
  \[
  q=
  \left(
  r_{\rm Li-A},
  s_{\rm pair},
  n_{\rm Li-solv},
  n_{\rm Li-lig},
  \omega_{\rm A},
  n_{\rm cluster},
  \rho_{\rm local},
  p_{\rm atm},
  \ldots
  \right).
  \]

  Every component of \(q\) must have:

  1. a unit,
  2. a physical meaning,
  3. a gradient,
  4. a mobility entry in \(D_x(q)\),
  5. a role in \(U_x(q)\), \(P_x(q)\), a basin \(A_i\), or a memory coordinate
  \(\psi_\alpha\).

  \section{Generator}

  \subsection{Free energy}

  Define molar free energy
  \[
  U_x:\Omega_x\to\R
  \qquad [\mathrm{J\,mol^{-1}}].
  \]

  It may be decomposed as
  \[
  U_x(q)
  =
  U_{\rm bonded}
  +
  U_{\rm rep}
  +
  U_{\rm disp}
  +
  U_{\rm Coul}
  +
  U_{\rm pol}
  +
  U_{\rm solv}
  +
  U_{\rm coord}
  +
  U_{\rm pack}
  +
  U_{\rm atm}
  +
  U_{\rm ext}.
  \]

  The exact decomposition is not important. What matters is that
  \[
  U_x(q)
  \]
  is a numerical function evaluable at quadrature points.

  \subsection{Thermal factor}

  Use the molar gas constant
  \[
  R=8.31446261815324\ \mathrm{J\,mol^{-1}\,K^{-1}}.
  \]

  Define
  \[
  \beta_{\rm mol}=\frac{1}{RT}.
  \]

  Then
  \[
  \exp[-\beta_{\rm mol}U_x(q)]
  \]
  is the Boltzmann factor.

  \subsection{Mobility tensor}

  Define
  \[
  D_x(q)\in\R^{d\times d}
  \qquad [\mathrm{m^2\,s^{-1}}].
  \]

  It must satisfy
  \[
  D_x(q)=D_x(q)^\top,
  \qquad
  D_x(q)\succeq0.
  \]

  If built from resistance:
  \[
  D_x(q)=k_BT\,R_x(q)^\#.
  \]

  Resistance may be assembled as
  \[
  R_x(q)
  =
  R_{\rm hydro}
  +
  R_{\rm crowd}
  +
  R_{\rm charge\ cloud}
  +
  R_{\rm atmosphere}
  +
  R_{\rm cage}
  +
  R_{\rm constraint}.
  \]

  \subsection{Reversible generator}

  For overdamped Smoluchowski dynamics:
  \[
  L_x f(q)
  =
  \frac{1}{\mu_x(q)}
  \nabla_q\cdot
  \left(
  \mu_x(q)D_x(q)\nabla_q f(q)
  \right).
  \]

  Equivalently:
  \[
  L_x f
  =
  D_x:\nabla_q\nabla_q f
  +
  \left[
  \nabla_q\cdot D_x
  -
  \beta_{\rm mol}D_x\nabla_q U_x
  \right]\cdot\nabla_q f.
  \]

  The equilibrium measure is
  \[
  \mu_x(\dd q)
  =
  Z_x^{-1}
  \exp[-\beta_{\rm mol}U_x(q)]
  \dd q.
  \]

  The model does not need to simulate \(L_x\) in production. It uses \(L_x\)
  through quadrature, committor/capacity solves, transition-path moment solves,
  and Mori inner products.

  \section{Charge polarization}

  Define formal-charge polarization:
  \[
  P_x(q)=\sum_{\alpha\in\mathcal Q} z_\alpha r_\alpha^{\rm unwrap}(q)
  \in\R^3.
  \]

  Here:

  \[
  \mathcal Q
  =
  \text{set of charge sites},
  \]
  \[
  z_\alpha
  =
  \text{dimensionless charge number of site }\alpha,
  \]
  \[
  r_\alpha^{\rm unwrap}(q)
  =
  \text{unwrapped position of charge site }\alpha.
  \]

  The gradient of charge polarization is
  \[
  G_P(q)=\nabla_q P_x(q)\in\R^{3\times d}.
  \]

  Component notation:
  \[
  (G_P)_{a\ell}(q)
  =
  \frac{\partial P_{x,a}(q)}{\partial q_\ell},
  \]
  where \(a\in\{x,y,z\}\).

  Using formal charge units, the final conductivity prefactor is
  \[
  \frac{F^2}{RT}.
  \]

  If physical charge units are used instead, the prefactor changes to
  \[
  \frac{1}{k_BT\,V}.
  \]

  This document uses formal charge units.

  \section{Projection basins}

  Choose \(N\) disjoint basins:
  \[
  \mathcal A=\{A_1,\ldots,A_N\},
  \]
  where
  \[
  A_i\subset\Omega_x,
  \qquad
  A_i\cap A_j=\varnothing \quad(i\ne j).
  \]

  Every coordinate \(q\) used in the finite model is assigned to exactly one
  basin:
  \[
  s(q)=i
  \quad
  \Longleftrightarrow
  \quad
  q\in A_i.
  \]

  Each basin \(A_i\) has a stoichiometry vector:
  \[
  \nu_i=(\nu_{i,a})_{a\in\mathcal C}.
  \]

  \[
  \nu_{i,a}
  =
  \text{number of conserved component }a\text{ in basin/motif }i.
  \]

  Examples:

  \[
  \nu_{i,\mathrm{Li}}=1
  \]
  for a motif containing one Li.

  \[
  \nu_{i,\mathrm{PF6}}=1
  \]
  for a motif containing one PF6.

  \[
  \nu_{i,\mathrm{FEC}}=1
  \]
  for a Li--FEC ligand motif.

  A basin is a coordinate set, not a terminal correction.

  \section{Quadrature}

  For each basin \(A_i\), choose quadrature points and weights:
  \[
  \{(q_{i\ell},w_{i\ell})\}_{\ell=1}^{n_i}.
  \]

  Each
  \[
  q_{i\ell}\in A_i.
  \]

  Weights satisfy
  \[
  \int_{A_i} f(q)\dd q
  \approx
  \sum_{\ell=1}^{n_i}w_{i\ell}f(q_{i\ell}).
  \]

  Weights may be generated by:

  1. tensor quadrature,
  2. sparse grid quadrature,
  3. finite element quadrature,
  4. Monte Carlo / importance quadrature,
  5. umbrella-sampling reweighting.

  The model only requires that the same weights are used consistently in all
  basin integrals.

  \section{Mass balance through chemical potentials}

  \subsection{Why chemical potentials are needed}

  Composition gives conserved totals \(C_a\). The projected basins \(A_i\) are
  motif states that consume conserved components according to \(\nu_{i,a}\).
  Therefore the basin populations \(c_i\) must satisfy:
  \[
  \sum_i\nu_{i,a}c_i=C_a
  \qquad
  \text{for every conserved component }a.
  \]

  This is achieved by component chemical potentials \(\lambda_a\).

  \subsection{Grand-canonical restricted basin weight}

  For each basin \(i\), define
  \[
  Z_i(\lambda)
  =
  \int_{A_i}
  \exp
  \left[
  -\beta_{\rm mol}U_x(q)
  +
  \sum_{a\in\mathcal C}\nu_{i,a}\lambda_a
  \right]\dd q.
  \]

  The dimensionless chemical potential \(\lambda_a\) is measured in units of
  \(RT\). It is not in J/mol. If using molar chemical potential
  \(\mu_a^{\rm chem}\), then
  \[
  \lambda_a=\frac{\mu_a^{\rm chem}}{RT}.
  \]

  Use standard concentration
  \[
  C^\circ=1000\ \mathrm{mol\,m^{-3}}.
  \]

  Then basin concentration is
  \[
  c_i(\lambda)=C^\circ Z_i(\lambda).
  \]

  Quadrature form:
  \[
  Z_i(\lambda)
  \approx
  \sum_{\ell=1}^{n_i}
  w_{i\ell}
  \exp
  \left[
  -\beta_{\rm mol}U_x(q_{i\ell})
  +
  \sum_{a\in\mathcal C}\nu_{i,a}\lambda_a
  \right].
  \]

  \subsection{Mass-balance residual}

  For every conserved component \(a\), define
  \[
  r_a(\lambda)
  =
  \sum_{i=1}^N\nu_{i,a}c_i(\lambda)-C_a.
  \]

  The chemical potential solve is:
  \[
  r_a(\lambda^\star)=0
  \qquad
  \forall a.
  \]

  \subsection{Jacobian}

  The derivative is
  \[
  \frac{\partial c_i}{\partial\lambda_b}
  =
  \nu_{i,b}c_i.
  \]

  Therefore
  \[
  J_{ab}(\lambda)
  =
  \frac{\partial r_a}{\partial\lambda_b}
  =
  \sum_{i=1}^N\nu_{i,a}\nu_{i,b}c_i(\lambda).
  \]

  \subsection{Newton algorithm}

  Input:

  \[
  C_a,\quad \nu_{i,a},\quad U_x,\quad A_i,\quad q_{i\ell},\quad w_{i\ell}.
  \]

  Tolerance:
  \[
  \tau_{\rm mass}=10^{-10}.
  \]

  Maximum iterations:
  \[
  N_{\rm Newton}=100.
  \]

  Initialization:
  \[
  \lambda_a^{(0)}
  =
  \log
  \left(
  \frac{\max(C_a,C_{\min})}{C^\circ}
  \right),
  \]
  with
  \[
  C_{\min}=10^{-30}\ \mathrm{mol\,m^{-3}}.
  \]

  Iteration \(k\):

  1. Compute all \(Z_i(\lambda^{(k)})\).

  2. Compute all
  \[
  c_i^{(k)}=C^\circ Z_i(\lambda^{(k)}).
  \]

  3. Compute residual
  \[
  r_a^{(k)}
  =
  \sum_i\nu_{i,a}c_i^{(k)}-C_a.
  \]

  4. Compute normalized residual
  \[
  \eta^{(k)}
  =
  \max_a
  \frac{|r_a^{(k)}|}
  {\max(C_a,C_{\min})}.
  \]

  5. Stop if
  \[
  \eta^{(k)}<\tau_{\rm mass}.
  \]

  6. Compute Jacobian
  \[
  J_{ab}^{(k)}
  =
  \sum_i\nu_{i,a}\nu_{i,b}c_i^{(k)}.
  \]

  7. Solve the linear Newton system
  \[
  J^{(k)}\Delta\lambda^{(k)}=-r^{(k)}.
  \]

  If \(J^{(k)}\) is singular or ill-conditioned, use Tikhonov-regularized solve:
  \[
  \left(J^{(k)}+\epsilon_J I\right)\Delta\lambda^{(k)}=-r^{(k)},
  \]
  where
  \[
  \epsilon_J=10^{-12}\max(1,\|J^{(k)}\|_2).
  \]

  8. Backtracking line search:

  Start with
  \[
  \alpha=1.
  \]

  Trial:
  \[
  \lambda_{\rm trial}=\lambda^{(k)}+\alpha\Delta\lambda^{(k)}.
  \]

  Compute
  \[
  \eta_{\rm trial}.
  \]

  Accept the step if
  \[
  \eta_{\rm trial}<\eta^{(k)}.
  \]

  If not accepted, set
  \[
  \alpha\leftarrow \frac12\alpha
  \]
  and try again.

  Fail if
  \[
  \alpha<2^{-40}.
  \]

  9. Update:
  \[
  \lambda^{(k+1)}=\lambda_{\rm trial}.
  \]

  After convergence:
  \[
  \lambda^\star=\lambda^{(k)}.
  \]

  The final concentration is:
  \[
  c_i=c_i(\lambda^\star).
  \]

  \subsection{Density quadrature weights}

  Define the concentration-density quadrature weight:
  \[
  W_{i\ell}
  =
  C^\circ
  w_{i\ell}
  \exp
  \left[
  -\beta_{\rm mol}U_x(q_{i\ell})
  +
  \sum_a\nu_{i,a}\lambda_a^\star
  \right].
  \]

  Then
  \[
  c_i=\sum_{\ell=1}^{n_i}W_{i\ell}.
  \]

  These \(W_{i\ell}\) are used in all basin averages.

  \section{Reactive fluxes}

  \subsection{Purpose}

  Reactive fluxes determine how often the projected process moves between basins.

  They produce:
  \[
  K_{ij},\quad Q_{ij}.
  \]

  \[
  K_{ij}
  \quad [\mathrm{mol\,m^{-3}\,s^{-1}}]
  \]
  is the symmetric concentration flux between basin \(i\) and basin \(j\).

  \[
  Q_{ij}
  \quad [\mathrm{s^{-1}}]
  \]
  is the finite-state transition rate from \(i\) to \(j\).

  \subsection{Committor definition}

  For an edge \((i,j)\), define the committor:
  \[
  q_{ij}(q)
  =
  \Pr_q(\tau_j<\tau_i).
  \]

  Here
  \[
  \tau_i=\inf\{t:X_t\in A_i\}.
  \]

  The committor satisfies:
  \[
  L_x q_{ij}=0
  \quad
  \text{on }\Omega_{ij}=\Omega_x\setminus(A_i\cup A_j),
  \]
  \[
  q_{ij}=0
  \quad
  \text{on }A_i,
  \]
  \[
  q_{ij}=1
  \quad
  \text{on }A_j.
  \]

  \subsection{How to solve the committor numerically}

  Choose finite element basis functions
  \[
  \{\varphi_m(q)\}_{m=1}^{M}
  \]
  on the transition domain \(\Omega_{ij}\).

  Approximate:
  \[
  q_{ij}(q)\approx \sum_{m=1}^{M}u_m\varphi_m(q).
  \]

  Boundary nodes on \(A_i\) are fixed to \(0\).
  Boundary nodes on \(A_j\) are fixed to \(1\).

  For interior degrees of freedom, solve:
  \[
  S u = f,
  \]
  where the stiffness matrix is
  \[
  S_{mn}
  =
  \int_{\Omega_{ij}}
  \nabla\varphi_m(q)^\top
  D_x(q)
  \nabla\varphi_n(q)
  \,\rho_x(\dd q),
  \]
  and \(\rho_x\) is the concentration measure:
  \[
  \rho_x(\dd q)
  =
  C^\circ
  \exp
  \left[
  -\beta_{\rm mol}U_x(q)
  +
  \sum_a\nu_{s(q),a}\lambda_a^\star
  \right]\dd q.
  \]

  The right-hand side \(f\) comes only from known boundary values after
  eliminating Dirichlet nodes.

  After solving \(u\), evaluate
  \[
  \nabla q_{ij}(q)
  \]
  at transition quadrature points.

  \subsection{Capacity flux}

  The symmetric capacity flux is:
  \[
  K_{ij}
  =
  \int_{\Omega_{ij}}
  \nabla q_{ij}(q)^\top
  D_x(q)
  \nabla q_{ij}(q)\,
  \rho_x(\dd q).
  \]

  Using transition quadrature points \(q_{ij,\ell}^{\Sigma}\), weights
  \(w_{ij,\ell}^{\Sigma}\), and committor gradients
  \[
  g_{ij,\ell}=\nabla q_{ij}(q_{ij,\ell}^{\Sigma}),
  \]
  compute:
  \[
  K_{ij}
  \approx
  \sum_{\ell}
  W_{ij,\ell}^{\Sigma}
  g_{ij,\ell}^{\top}
  D_x(q_{ij,\ell}^{\Sigma})
  g_{ij,\ell}.
  \]

  The transition-surface concentration weight is:
  \[
  W_{ij,\ell}^{\Sigma}
  =
  C^\circ
  w_{ij,\ell}^{\Sigma}
  \exp
  \left[
  -\beta_{\rm mol}U_x(q_{ij,\ell}^{\Sigma})
  +
  \sum_a\nu_{s(q_{ij,\ell}^{\Sigma}),a}\lambda_a^\star
  \right].
  \]

  Set:
  \[
  K_{ji}=K_{ij}.
  \]

  Set:
  \[
  K_{ii}=0.
  \]

  \subsection{Finite generator}

  For \(i\ne j\):
  \[
  Q_{ij}=\frac{K_{ij}}{c_i}.
  \]

  For diagonal entries:
  \[
  Q_{ii}=-\sum_{j\ne i}Q_{ij}.
  \]

  Validation:

  1. Row sums:
  \[
  \sum_j Q_{ij}=0.
  \]

  2. Detailed balance:
  \[
  c_iQ_{ij}=c_jQ_{ji}.
  \]

  3. Stationarity:
  \[
  \sum_i c_iQ_{ij}=0.
  \]

  \section{Transition-path displacement moments}

  \subsection{Purpose}

  Transition-path moments determine charge transported during state changes.

  They produce:
  \[
  d_{ij},\quad M_{ij}.
  \]

  \[
  d_{ij}\in\R^3
  \quad [\mathrm{m}]
  \]
  is mean formal-charge displacement for transition \(i\to j\).

  \[
  M_{ij}\in\R^{3\times3}
  \quad [\mathrm{m^2}]
  \]
  is second moment of formal-charge displacement for transition \(i\to j\).

  \subsection{Reactive path sampling algorithm}

  For edge \((i,j)\), use the committor \(q_{ij}\).

  The process conditioned to reach \(A_j\) before \(A_i\) has Doob-transformed
  drift:
  \[
  b^{(ij)}(q)
  =
  b(q)+2D_x(q)\nabla\log q_{ij}(q),
  \]
  where
  \[
  b(q)=\nabla_q\cdot D_x(q)-\beta_{\rm mol}D_x(q)\nabla_q U_x(q)
  \]
  is the original overdamped drift.

  The conditioned short-path SDE is:
  \[
  \dd q_t
  =
  b^{(ij)}(q_t)\dd t
  +
  \sqrt{2D_x(q_t)}\dd W_t.
  \]

  Initial points are sampled from the reactive exit distribution on the
  \(i\)-side boundary:
  \[
  p_{\rm exit}(q)
  \propto
  \rho_x(q)
  \,n_{ij}(q)^\top D_x(q)\nabla q_{ij}(q),
  \]
  where \(n_{ij}\) is the outward normal from \(A_i\) toward \(A_j\).

  For each sampled path:

  1. Start at \(q_0\) from \(p_{\rm exit}\).

  2. Integrate the conditioned SDE until hitting \(A_j\).

  3. Record endpoint \(q_\tau\).

  4. Compute charge displacement:
  \[
  \Delta P_{ij,\ell}=P_x(q_\tau)-P_x(q_0).
  \]

  This is a local transition-path calculation, not a long-time GK integral.

  \subsection{Weighted moment estimator}

  Given path samples
  \[
  \{(\Delta P_{ij,\ell},w_{ij,\ell}^{\rm path})\}_{\ell=1}^{n_{ij}^{\rm path}},
  \]
  compute:
  \[
  S_{ij}^{\rm path}=\sum_\ell w_{ij,\ell}^{\rm path}.
  \]

  Then:
  \[
  d_{ij}
  =
  \frac{1}{S_{ij}^{\rm path}}
  \sum_\ell
  w_{ij,\ell}^{\rm path}\Delta P_{ij,\ell}.
  \]

  \[
  M_{ij}
  =
  \frac{1}{S_{ij}^{\rm path}}
  \sum_\ell
  w_{ij,\ell}^{\rm path}
  \Delta P_{ij,\ell}\Delta P_{ij,\ell}^{\top}.
  \]

  Set reverse moments:
  \[
  d_{ji}=-d_{ij},
  \]
  \[
  M_{ji}=M_{ij}.
  \]

  Set:
  \[
  d_{ii}=0,
  \qquad
  M_{ii}=0.
  \]

  \subsection{Displacement-scale validation}

  For every transition displacement sample, require:
  \[
  \|\Delta P_{ij,\ell}\| \le L_{\max}.
  \]

  For molecular local transitions, choose for example:
  \[
  L_{\max}=10^{-8}\ \mathrm{m}.
  \]

  A transition displacement of order meters must fail, because \(P_x\) is an
  unwrapped molecular charge displacement in meters.

  \section{Within-state self-current}

  \subsection{Purpose}

  Self-current captures continuous unbounded charge mobility while the projected
  state does not change.

  It produces:
  \[
  D_i^{\rm self}.
  \]

  \[
  D_i^{\rm self}\in\R^{3\times3}
  \qquad [\mathrm{m^2\,s^{-1}}].
  \]

  \subsection{Unbounded transport projector}

  For each basin \(A_i\), define:
  \[
  \Pi_i(q)\in\R^{d\times d}.
  \]

  \(\Pi_i(q)\) projects coordinate motion onto directions that produce
  unbounded DC charge transport.

  Examples:

  1. Free ion translation:
  \[
  \Pi_i=I
  \]
  on that ion's translational coordinates.

  2. Bound internal pair vibration:
  \[
  \Pi_i=0
  \]
  for the bounded internal coordinate.

  3. Neutral pair center-of-mass translation:
  \[
  \Pi_i=I_{\rm COM}
  \]
  but
  \[
  G_P\Pi_i=0
  \]
  if the neutral pair is exactly charge-neutral and co-moving.

  4. Partner-switch or identity-diffusion coordinate:
  \[
  \Pi_i
  \]
  is nonzero only along the unbounded identity coordinate.

  \subsection{Self-current integral}

  For each state:
  \[
  D_i^{\rm self}
  =
  \frac{1}{c_i}
  \int_{A_i}
  G_P(q)
  \Pi_i(q)
  D_x(q)
  \Pi_i(q)^\top
  G_P(q)^\top
  \,\rho_x(\dd q).
  \]

  Quadrature:
  \[
  D_i^{\rm self}
  \approx
  \frac{1}{c_i}
  \sum_{\ell=1}^{n_i}
  W_{i\ell}
  G_P(q_{i\ell})
  \Pi_i(q_{i\ell})
  D_x(q_{i\ell})
  \Pi_i(q_{i\ell})^\top
  G_P(q_{i\ell})^\top.
  \]

  Validation:

  \[
  D_i^{\rm self}=D_i^{\rm self\top},
  \qquad
  D_i^{\rm self}\succeq0.
  \]

  \subsection{Multi-center charge covariance}

  If a basin contains multiple charged centers, do not use
  \[
  z_{\rm net}^2D_{\rm COM}.
  \]

  Let
  \[
  z_i=(z_1,\ldots,z_m)^\top
  \]
  be the charge vector of centers in basin \(i\), and let
  \[
  D_i^{\rm centers}\in\R^{m\times m}
  \]
  be the state-local mobility covariance of those centers along one Cartesian
  axis.

  Then charge mobility is:
  \[
  D_Q(i)=z_i^\top D_i^{\rm centers}z_i.
  \]

  For a Li--anion pair:
  \[
  z=(+1,-1)^\top.
  \]

  Thus:
  \[
  D_Q
  =
  D_{\rm Li}
  +
  D_{\rm A}
  -
  2D_{\rm Li,A}.
  \]

  If
  \[
  D_{\rm Li,A}<0,
  \]
  then the cross term raises conductivity.

  \section{Direct projected diffusivity-density}

  Define:
  \[
  B
  =
  \sum_i c_iD_i^{\rm self}
  +
  \frac{1}{2}\sum_i\sum_jK_{ij}M_{ij}.
  \]

  Units:
  \[
  [c_iD_i^{\rm self}]
  =
  \mathrm{mol\,m^{-3}}\cdot\mathrm{m^2\,s^{-1}}
  =
  \mathrm{mol\,m^{-1}\,s^{-1}}.
  \]

  \[
  [K_{ij}M_{ij}]
  =
  \mathrm{mol\,m^{-3}\,s^{-1}}\cdot \mathrm{m^2}
  =
  \mathrm{mol\,m^{-1}\,s^{-1}}.
  \]

  So:
  \[
  [B]=\mathrm{mol\,m^{-1}\,s^{-1}}.
  \]

  \section{Finite-state memory correction}

  \subsection{State drift}

  For each state:
  \[
  b_i=\sum_jQ_{ij}d_{ij}\in\R^3.
  \]

  Component:
  \[
  b_{i,a}=\sum_jQ_{ij}d_{ij,a}.
  \]

  \subsection{Poisson equation}

  For each Cartesian axis \(a\), solve:
  \[
  -Q\chi_a=b_a.
  \]

  Because \(Q\) has a zero eigenvalue, impose gauge:
  \[
  \sum_i c_i\chi_{i,a}=0.
  \]

  \subsection{Numerical solve}

  Construct block matrix:
  \[
  \mathcal M
  =
  \begin{bmatrix}
  -Q & \one\\
  c^\top & 0
  \end{bmatrix}.
  \]

  Construct right-hand side:
  \[
  y_a=
  \begin{bmatrix}
  b_a\\
  0
  \end{bmatrix}.
  \]

  Solve:
  \[
  \mathcal M
  \begin{bmatrix}
  \chi_a\\
  \lambda_a^{Q}
  \end{bmatrix}
  =
  y_a.
  \]

  Here \(\lambda_a^{Q}\) is a Lagrange multiplier enforcing the gauge.

  The solution vector is:
  \[
  \chi_a=(\chi_{1,a},\ldots,\chi_{N,a})^\top.
  \]

  \subsection{Correction tensor}

  \[
  C^Q_{ab}
  =
  \sum_i c_i b_{i,a}\chi_{i,b}.
  \]

  Thus:
  \[
  C^Q\in\R^{3\times3}.
  \]

  Symmetrize numerically:
  \[
  C^Q\leftarrow\frac12(C^Q+C^{Q\top}).
  \]

  \section{Continuous Mori memory correction}

  \subsection{Memory coordinates}

  Choose memory coordinates:
  \[
  \Psi(q)
  =
  (\psi_1(q),\ldots,\psi_m(q))^\top.
  \]

  Examples:

  \[
  \psi_{\rm atm}(q)
  =
  \text{ionic-atmosphere polarization coordinate},
  \]
  \[
  \psi_{\rm cage}(q)
  =
  \text{solvent-cage/backjump coordinate},
  \]
  \[
  \psi_{\rm partner}(q)
  =
  \text{partner residence coordinate},
  \]
  \[
  \psi_{\rm orient}(q)
  =
  \text{anion orientation coordinate}.
  \]

  \subsection{State-indicator orthogonalization}

  For each state \(A_i\):
  \[
  \bar\psi_i
  =
  \frac{1}{c_i}
  \int_{A_i}\Psi(q)\rho_x(\dd q).
  \]

  Quadrature:
  \[
  \bar\psi_i
  =
  \frac{1}{c_i}
  \sum_{\ell}W_{i\ell}\Psi(q_{i\ell}).
  \]

  Define:
  \[
  \Psi^\perp(q)=\Psi(q)-\bar\psi_{s(q)}.
  \]

  This prevents double-counting memory already represented by finite state
  indicators.

  If \(\bar\psi_{s(q)}\) is piecewise constant, then inside a basin:
  \[
  \nabla_q\Psi^\perp(q)=\nabla_q\Psi(q).
  \]

  At basin boundaries, discontinuities are already handled by the finite-state
  Poisson correction.

  \subsection{Gradient}

  Define:
  \[
  G_\psi(q)
  =
  \nabla_q\Psi^\perp(q)
  \in\R^{m\times d}.
  \]

  Component:
  \[
  (G_\psi)_{\alpha\ell}
  =
  \frac{\partial\psi_\alpha^\perp}{\partial q_\ell}.
  \]

  \subsection{Mori matrix}

  \[
  A_{\alpha\gamma}
  =
  \int
  \nabla\psi_\alpha^\perp(q)^\top
  D_x(q)
  \nabla\psi_\gamma^\perp(q)
  \,\rho_x(\dd q).
  \]

  Quadrature:
  \[
  A
  \approx
  \sum_i\sum_{\ell}
  W_{i\ell}
  G_\psi(q_{i\ell})
  D_x(q_{i\ell})
  G_\psi(q_{i\ell})^\top.
  \]

  Units:
  \[
  [A]=\mathrm{mol\,m^{-3}}\cdot \mathrm{m^2\,s^{-1}}\cdot [\psi]^2/[q]^2.
  \]

  \subsection{Current coupling}

  \[
  h_{\alpha a}
  =
  \int
  \nabla\psi_\alpha^\perp(q)^\top
  D_x(q)
  \nabla P_a(q)
  \,\rho_x(\dd q).
  \]

  Quadrature:
  \[
  h
  \approx
  \sum_i\sum_\ell
  W_{i\ell}
  G_\psi(q_{i\ell})
  D_x(q_{i\ell})
  G_P(q_{i\ell})^\top.
  \]

  \[
  h\in\R^{m\times3}.
  \]

  \subsection{Mori correction}

  Compute pseudoinverse:
  \[
  A^\#
  \]
  by eigendecomposition as defined in Section 1.

  Then:
  \[
  C^\psi=h^\top A^\# h.
  \]

  \[
  C^\psi\in\R^{3\times3}.
  \]

  Symmetrize numerically:
  \[
  C^\psi\leftarrow\frac12(C^\psi+C^{\psi\top}).
  \]

  \section{Projected diffusivity-density}

  \[
  D_{\rm proj}
  =
  B-C^Q-C^\psi.
  \]

  Symmetrize:
  \[
  D_{\rm proj}\leftarrow\frac12(D_{\rm proj}+D_{\rm proj}^{\top}).
  \]

  If small negative eigenvalues occur from numerical roundoff, project to PSD:

  1. Eigendecompose:
  \[
  D_{\rm proj}=V\diag(\eta_1,\eta_2,\eta_3)V^\top.
  \]

  2. Set:
  \[
  \eta_r^+=\max(\eta_r,0).
  \]

  3. Reconstruct:
  \[
  D_{\rm proj}^{+}=V\diag(\eta_1^+,\eta_2^+,\eta_3^+)V^\top.
  \]

  Use this projection only if:
  \[
  \min_r\eta_r\ge -\tau_D\max_r|\eta_r|,
  \]
  for example
  \[
  \tau_D=10^{-10}.
  \]

  If the negative eigenvalue is larger than tolerance, fail because the primitive
  set is physically inconsistent.

  \section{Final scalar conductivity}

  The isotropic projected diffusivity-density is:
  \[
  D_{\rm iso}
  =
  \frac13\Tr(D_{\rm proj}).
  \]

  Conductivity:
  \[
  \sigma_{\rm S/m}
  =
  \frac{F^2}{RT}D_{\rm iso}.
  \]

  Here:
  \[
  F=96485.33212\ \mathrm{C\,mol^{-1}}.
  \]

  Units:
  \[
  \frac{F^2}{RT}
  \cdot
  \mathrm{mol\,m^{-1}\,s^{-1}}
  =
  \mathrm{S\,m^{-1}}.
  \]

  Convert:
  \[
  \sigma_{\rm mS/cm}=10\sigma_{\rm S/m}.
  \]

  Final result:
  \[
  \boxed{
  \sigma_{\rm mS/cm}
  =
  10\frac{F^2}{3RT}
  \Tr
  \left[
  \sum_i c_iD_i^{\rm self}
  +
  \frac12\sum_{i,j}K_{ij}M_{ij}
  -
  C^Q
  -
  h^\top A^\#h
  \right].
  }
  \]

  \section{Complete numerical algorithm}

  Input:
  \[
  x=(T,\{C_a\},\{\phi_s\},\mathcal P).
  \]

  Output:
  \[
  \sigma_{\rm mS/cm}(x).
  \]

  Algorithm:

  \begin{enumerate}

  \item Build coordinate space \(\Omega_x\) and coordinate vector \(q\).

  \item Build numerical functions:
  \[
  U_x(q),
  \qquad
  D_x(q),
  \qquad
  P_x(q),
  \qquad
  G_P(q)=\nabla P_x(q).
  \]

  \item Build basin set:
  \[
  A_1,\ldots,A_N.
  \]

  \item Assign stoichiometry:
  \[
  \nu_{i,a}.
  \]

  \item Generate quadrature:
  \[
  (q_{i\ell},w_{i\ell})
  \quad
  \text{for each basin }A_i.
  \]

  \item Solve chemical potentials \(\lambda^\star\) by damped Newton:

  \[
  r_a(\lambda)=\sum_i\nu_{i,a}C^\circ Z_i(\lambda)-C_a.
  \]

  \[
  J_{ab}(\lambda)=\sum_i\nu_{i,a}\nu_{i,b}c_i(\lambda).
  \]

  \[
  J\Delta\lambda=-r.
  \]

  Backtrack until normalized residual decreases. Stop when
  \[
  \max_a\frac{|r_a|}{\max(C_a,C_{\min})}<10^{-10}.
  \]

  \item Compute density weights:
  \[
  W_{i\ell}
  =
  C^\circ
  w_{i\ell}
  \exp
  \left[
  -\frac{U_x(q_{i\ell})}{RT}
  +
  \sum_a\nu_{i,a}\lambda_a^\star
  \right].
  \]

  \item Compute populations:
  \[
  c_i=\sum_\ell W_{i\ell}.
  \]

  \item For every transition edge \((i,j)\), solve the committor PDE:
  \[
  L_xq_{ij}=0,
  \quad
  q_{ij}|_{A_i}=0,
  \quad
  q_{ij}|_{A_j}=1.
  \]

  Use finite elements or finite volumes. The linear system is:
  \[
  S u=f,
  \]
  with
  \[
  S_{mn}
  =
  \int
  \nabla\varphi_m^\top D_x\nabla\varphi_n
  \,\rho_x(\dd q).
  \]

  \item Compute capacity flux:
  \[
  K_{ij}
  =
  \int
  \nabla q_{ij}^\top D_x\nabla q_{ij}
  \,\rho_x(\dd q).
  \]

  Set:
  \[
  K_{ji}=K_{ij}.
  \]

  \item Compute generator:
  \[
  Q_{ij}=K_{ij}/c_i,
  \qquad
  Q_{ii}=-\sum_{j\ne i}Q_{ij}.
  \]

  \item For every transition edge, sample local conditioned transition paths with
  Doob drift:
  \[
  b^{(ij)}(q)
  =
  \nabla\cdot D_x(q)-\frac{1}{RT}D_x(q)\nabla U_x(q)
  +
  2D_x(q)\nabla\log q_{ij}(q).
  \]

  For each path, compute:
  \[
  \Delta P=P(q_{\rm end})-P(q_{\rm start}).
  \]

  \item Compute transition moments:
  \[
  d_{ij}
  =
  \frac{\sum_\ell w_\ell^{\rm path}\Delta P_\ell}
  {\sum_\ell w_\ell^{\rm path}},
  \]
  \[
  M_{ij}
  =
  \frac{\sum_\ell w_\ell^{\rm path}\Delta P_\ell\Delta P_\ell^\top}
  {\sum_\ell w_\ell^{\rm path}}.
  \]

  Set:
  \[
  d_{ji}=-d_{ij},
  \qquad
  M_{ji}=M_{ij}.
  \]

  \item Define unbounded projectors:
  \[
  \Pi_i(q).
  \]

  \item Compute self-current tensors:
  \[
  D_i^{\rm self}
  =
  \frac{1}{c_i}
  \sum_\ell
  W_{i\ell}
  G_P(q_{i\ell})
  \Pi_i(q_{i\ell})
  D_x(q_{i\ell})
  \Pi_i(q_{i\ell})^\top
  G_P(q_{i\ell})^\top.
  \]

  \item Compute direct tensor:
  \[
  B=
  \sum_i c_iD_i^{\rm self}
  +
  \frac12\sum_{i,j}K_{ij}M_{ij}.
  \]

  \item Compute drift:
  \[
  b_i=\sum_jQ_{ij}d_{ij}.
  \]

  \item For each Cartesian axis \(a\), solve:
  \[
  \begin{bmatrix}
  -Q & \one\\
  c^\top & 0
  \end{bmatrix}
  \begin{bmatrix}
  \chi_a\\
  \lambda_a^Q
  \end{bmatrix}
  =
  \begin{bmatrix}
  b_a\\
  0
  \end{bmatrix}.
  \]

  \item Compute finite-state correction:
  \[
  C^Q_{ab}=\sum_i c_i b_{i,a}\chi_{i,b}.
  \]

  \item Build memory coordinates:
  \[
  \Psi(q).
  \]

  \item Orthogonalize memory coordinates by state:
  \[
  \Psi^\perp(q)=\Psi(q)-\bar\Psi_{s(q)}.
  \]

  \item Compute memory gradient:
  \[
  G_\psi(q)=\nabla\Psi^\perp(q).
  \]

  \item Compute Mori matrix:
  \[
  A=
  \sum_i\sum_\ell
  W_{i\ell}
  G_\psi(q_{i\ell})
  D_x(q_{i\ell})
  G_\psi(q_{i\ell})^\top.
  \]

  \item Compute current coupling:
  \[
  h=
  \sum_i\sum_\ell
  W_{i\ell}
  G_\psi(q_{i\ell})
  D_x(q_{i\ell})
  G_P(q_{i\ell})^\top.
  \]

  \item Compute pseudoinverse \(A^\#\) by eigendecomposition.

  \item Compute continuous Mori correction:
  \[
  C^\psi=h^\top A^\#h.
  \]

  \item Compute projected diffusivity-density:
  \[
  D_{\rm proj}=B-C^Q-C^\psi.
  \]

  \item Compute scalar conductivity:
  \[
  \sigma_{\rm mS/cm}
  =
  10\frac{F^2}{3RT}\Tr(D_{\rm proj}).
  \]

  \end{enumerate}

  \section{Effect identification}

  An effect is any composition, coordinate, or parameter perturbation \(\theta\)
  that changes at least one primitive:
  \[
  c,\ K,\ Q,\ d,\ M,\ D^{\rm self},\ A,\ h.
  \]

  To identify the effect numerically, perturb \(\theta\) by a small amount
  \[
  \delta\theta.
  \]

  Compute central differences:
  \[
  \partial_\theta X
  \approx
  \frac{X(\theta+\delta\theta)-X(\theta-\delta\theta)}
  {2\delta\theta}
  \]
  for
  \[
  X\in
  \{c,K,d,M,D^{\rm self},A,h,\sigma\}.
  \]

  Primitive ownership scores:

  \[
  S_c(\theta)=\|\partial_\theta c\|_2,
  \]

  \[
  S_K(\theta)=\|\partial_\theta K\|_F,
  \]

  \[
  S_{dM}(\theta)=
  \|\partial_\theta d\|_F+\|\partial_\theta M\|_F,
  \]

  \[
  S_{D}(\theta)=\|\partial_\theta D^{\rm self}\|_F,
  \]

  \[
  S_{Ah}(\theta)=\|\partial_\theta A\|_F+\|\partial_\theta h\|_F.
  \]

  Assign the effect to the largest nonzero ownership score.

  \section{Conductivity-effect ledger}

  \[
  \begin{array}{lll}
  \toprule
  \text{Effect} & \text{Primitive location} & \text{Computation}\\
  \midrule
  \text{Free-ion fraction} & c_i & \partial c_i/\partial\theta\\
  \text{Ion association} & c_i,K_{ij} & \partial c,\partial K\\
  \text{SSIP/CIP balance} & c_i,K_{ij},D_i^{\rm self} & \partial c,\partial K,\partial D\\
  \text{Aggregate formation} & c_i,K_{ij},M_{ij},D_i^{\rm self} & \partial c,\partial K,\partial M,\partial D\\
  \text{Neutral ligand coordination} & c_i,K_{ij},A,h & \partial c,\partial K,\partial A,\partial h\\
  \text{Solvation competition} & U_x\to c_i,K_{ij} & \partial U\to\partial c,\partial K\\
  \text{Coulomb association} & U_x\to c_i,K_{ij} & \partial U\to\partial c,\partial K\\
  \text{Born/desolvation} & U_x\to c_i,K_{ij} & \partial U\to\partial c,\partial K\\
  \text{Activity/nonideality} & U_x\to c_i,K_{ij} & \partial U\to\partial c,\partial K\\
  \text{Viscosity} & D_x\to D_i^{\rm self},K_{ij},A & \partial D\to\partial D^{\rm self},\partial K,\partial A\\
  \text{Hydrodynamic size} & D_x\to D_i^{\rm self},K_{ij} & \partial D\\
  \text{Anion shape/denticity} & U_x,D_x,K_{ij},A,h & \partial U,\partial D,\partial K,\partial A,\partial h\\
  \text{Charge-cloud radius} & U_x,D_x,A,h & \partial U,\partial D,\partial A,\partial h\\
  \text{Local crowding/free volume} & U_x,D_x,K_{ij},A,h & \partial U,\partial D,\partial K,\partial A,\partial h\\
  \text{Li--anion co-motion} & D_i^{\rm self},M_{ij} & D_{\rm Li,A}>0\\
  \text{Li--anion anti-correlation} & D_i^{\rm self},M_{ij} & D_{\rm Li,A}<0\\
  \text{Partner switching} & K_{ij},d_{ij},M_{ij},A,h & \partial K,\partial d,\partial M,\partial A,\partial h\\
  \text{Identity diffusion} & K_{ij},d_{ij},M_{ij} & \partial K,\partial d,\partial M\\
  \text{Structural hopping} & K_{ij},d_{ij},M_{ij} & \partial K,\partial d,\partial M\\
  \text{Bound neutral polarization} & A,h & \partial A,\partial h\\
  \text{Cage/backjump} & C^Q \text{ or } A,h & \partial C^Q,\partial A,\partial h\\
  \text{Ion-atmosphere relaxation} & A,h & \partial A,\partial h\\
  \text{Electrophoretic drag} & D_x \text{ or } A,h & \partial D,\partial A,\partial h\\
  \text{Debye screening} & U_x,D_x,A,h & \partial U,\partial D,\partial A,\partial h\\
  \text{Solvent-shell residence} & c_i,K_{ij},A,h & \partial c,\partial K,\partial A,\partial h\\
  \text{Ligand-shell residence} & c_i,K_{ij},A,h & \partial c,\partial K,\partial A,\partial h\\
  \text{Temperature} & RT,U_x,D_x,K_{ij},F^2/(RT) & \partial_T\sigma\\
  \text{Dielectric response} & U_x,D_x,A,h & \partial U,\partial D,\partial A,\partial h\\
  \text{Salt concentration} & c_i,U_x,D_x,K_{ij},A,h & \partial c,\partial U,\partial D,\partial K,\partial A,\partial h\\
  \text{Additive microviscosity} & D_x,K_{ij} & \partial D,\partial K\\
  \text{Bulky-anion/ligand conductivity increase} & D_i^{\rm self}\text{ or }M_{ij} &
  D_{\rm Li,A|i}<0\text{ or partner-switch }M_{ij}\\
  \bottomrule
  \end{array}
  \]

  \section{Basis refinement without production Green--Kubo}

  Let the current memory basis be
  \[
  \Phi=(\phi_1,\ldots,\phi_m).
  \]

  For a candidate missing coordinate \(g(q)\), compute:

  \[
  A_{gg}
  =
  \int
  \nabla g^\top D_x\nabla g\,\rho_x(\dd q),
  \]

  \[
  A_{g\Phi}
  =
  \int
  \nabla g^\top D_x\nabla\Phi\,\rho_x(\dd q),
  \]

  \[
  h_g
  =
  \int
  \nabla g^\top D_xG_P^\top\,\rho_x(\dd q).
  \]

  Let
  \[
  A_{\Phi\Phi}
  =
  \int
  \nabla\Phi^\top D_x\nabla\Phi\,\rho_x(\dd q),
  \]
  \[
  h_\Phi
  =
  \int
  \nabla\Phi^\top D_xG_P^\top\,\rho_x(\dd q).
  \]

  Remove what is already represented:
  \[
  \tilde h_g
  =
  h_g-A_{g\Phi}A_{\Phi\Phi}^\#h_\Phi.
  \]

  \[
  \tilde A_{gg}
  =
  A_{gg}-A_{g\Phi}A_{\Phi\Phi}^\#A_{\Phi g}.
  \]

  Score:
  \[
  S(g)
  =
  \frac{\|\tilde h_g\|_2^2}{\tilde A_{gg}}.
  \]

  Add:
  \[
  g^\star=\argmax_g S(g).
  \]

  Stop when:
  \[
  \max_gS(g)<\tau_{\rm basis}
  \]
  and
  \[
  |\sigma_{n+1}-\sigma_n|<\tau_\sigma.
  \]

  This is the production convergence procedure. It uses the generator integrals,
  not a long-time GK integral.

  \section{Validation rules for implementation}

  A numerical implementation must fail if:

  \[
  c_i\le0,
  \]
  \[
  K_{ij}<0,
  \]
  \[
  K_{ij}\ne K_{ji},
  \]
  \[
  \sum_jQ_{ij}\ne0,
  \]
  \[
  c_iQ_{ij}\ne c_jQ_{ji},
  \]
  \[
  d_{ji}\ne-d_{ij},
  \]
  \[
  M_{ji}\ne M_{ij},
  \]
  \[
  D_i^{\rm self}\not\succeq0,
  \]
  \[
  A\not\succeq0,
  \]
  or if
  \[
  \|\Delta P_{ij,\ell}\|>L_{\max}.
  \]

  The projected analytical estimator should consume numerical generator inputs or
  primitive arrays. A property database row with only recipe composition and
  empirical conductivity should not be converted into a fake generator. Existing
  recipe-derived validators that call recipe evaluators are descriptor-surrogate
  validators, not the full projected analytical model. :contentReference[oaicite:1]{index=1}

  \section{Final implementation formula}

  \[
  \boxed{
  \sigma_{\rm mS/cm}(x)
  =
  10\frac{F^2}{3RT}
  \Tr
  \left[
  \sum_i c_iD_i^{\rm self}
  +
  \frac12\sum_{i,j}K_{ij}M_{ij}
  -
  C^Q
  -
  h^\top A^\#h
  \right].
  }
  \]

  Every symbol in this formula is computed by the explicit steps above.

  \end{document}























  
  The list below is the conductivity-effect checklist I would use.

  ---

  ## A. Readout and current-correlation effects

  | Effect tests must identify                     | What the test should verify
  |
  | ---------------------------------------------- |
  ---------------------------------------------------------------------------------------------------------------------------------------
  ---------------------------------------------------------------------------------------------------------------------------------------
  ---------------------------------------------------------------------------------------------- |
  | **Nernst–Einstein/direct carrier term**        | One charged diffusive state gives (D_Q=z^2D) and nonzero σ.
  |
  | **Poisson/Mori current-memory corrector**      | A correlated two-state jump process gives direct minus corrector, not direct-only σ.
  |
  | **Green–Kubo cross-current covariance**        | For opposite charges, negative Li–anion velocity covariance increases σ.
  |
  | **State-current drift coupling**               | Nonzero state-conditioned drift produces nonzero corrector; zero drift gives zero
  corrector only when physically justified.
  |
  | **Finite-basis closure**                       | Projected σ must be compared to trajectory GK/EH; missing variables are selected by
  unresolved current residual, not intuition. Candidate trajectory observables include Li–anion separation, Li–O coordination, CIP/SSIP
  label, solvent-shell identity, aggregate size, anion orientation/denticity, local free volume, and previous-hop coordinates only when
  persistent.  |
  | **Null current mode / divergent conductivity** | If current couples to a null energy mode, the model must fail loudly, not pseudo-
  invert away the problem.
  |
  | **Pure motif exchange**                        | Zero-displacement motif exchange must produce no direct conductivity.
  |
  | **Structural hop**                             | A state-changing event contributes to σ only if it has a physically defined charge-
  displacement moment and reversible partner event.
  |
  | **Event reversibility**                        | Every nonzero displacement event must have reverse flux with opposite displacement.
  |
  | **Detailed balance / stationarity**            | Finite σ is invalid if row conservation, stationarity, or detailed balance fail.
  |

  ---

  ## B. Speciation and population effects

  | Effect tests must identify           | What the test should verify
  |
  | ------------------------------------ |
  ---------------------------------------------------------------------------------------------------------------------------------------
  --------------------------------------------------------------------------- |
  | **Free-ion fraction**                | Increasing association strength lowers free carrier fraction.
  |
  | **Shared-Li mass balance**           | Mixed salts must share one Li pool; Li cannot be counted once per salt source. The plan
  explicitly requires shared Li, per-anion free concentrations, SSIP, CIP, Li-additive, and aggregate motif concentrations.  |
  | **SSIP population**                  | Solvent-separated ion-pair concentration changes (c_i), (D_i), (M_{ij}), and memory modes.
  |
  | **CIP population**                   | Contact pairs suppress or modify charge mobility through state-local multi-center covariance,
  not scalar net charge.                                                                                               |
  | **Li-additive coordination**         | FEC/VC/other neutral ligands must create Li-ligand motif population when coordination
  descriptors imply it. A test must fail if FEC only changes viscosity and never creates an FEC-coordinated motif.             |
  | **Li–additive–anion ternary motifs** | A Li-ligand shell near an anion must be distinguishable from free Li + neutral additive
  background.                                                                                                                |
  | **Aggregate formation**              | Li(_2)A(^+), LiA(_2^-), Li(_2)A(_2), and higher clusters must appear when stoichiometry and
  free energies require them.                                                                                            |
  | **Mixed-salt clusters**              | Generic stoichiometric clusters must be possible; no LiPF6/LiFSI-specific shortcut.
  |
  | **Neutral clusters**                 | Neutral multi-charge motifs must not be inserted as independent charged carriers, but they may
  carry internal charge mobility through (z^T D z).                                                                   |
  | **Multivalent ions**                 | Charge-balanced multivalent recipes must build the required neutral/charged clusters.
  |
  | **Activity / nonideality**           | Ionic-strength and hard-sphere activity effects must alter populations, not final σ.
  |
  | **Coulomb association**              | Scaling Coulomb interaction must change pair/cluster speciation.
  |
  | **Born desolvation**                 | Cavity/Born radius changes pair and cluster free energies.
  |
  | **Coordination free energy**         | Donor/acceptor/coordination affinity changes Li binding, SSIP/CIP split, and additive motifs.
  |
  | **Steric cluster penalty**           | Molecular volume and hard-sphere packing penalize clusters and change concentrations.
  |
  | **Cluster entropy penalty**          | Larger clusters must pay configurational/association entropy unless explicitly compensated.
  |

  The generic speciation module enumerates contact pairs, solvent-separated pairs, charged triplets, neutral clusters, and higher charged
  clusters, and solves mass balance over arbitrary charged components rather than a lithium-only algebra.

  ---

  ## C. State-local mobility and obstruction effects

  | Effect tests must identify               | What the test should verify
  |
  | ---------------------------------------- |
  ---------------------------------------------------------------------------------------------------------------------------------------
  -------------------------------------------------------------------------------- |
  | **Viscosity**                            | Higher mixture viscosity lowers local diffusivity and transition rates.
  |
  | **Temperature**                          | (k_BT), viscosity, rates, and (F^2/RT) all respond consistently.
  |
  | **Hydrodynamic radius**                  | Larger carrier or cluster radius lowers mobility.
  |
  | **Cluster hydrodynamic radius**          | Cluster diffusivity must use cluster geometry/radius, not free-ion radius.
  |
  | **Shape friction**                       | Bulky/asymmetric anions must change mobility via shape/ligand-field descriptors, not salt
  names. Hostile tests explicitly require bulky/asymmetric anion rows to fail if no structure-derived shape factor is present.  |
  | **Charge-cloud radius / delocalization** | Charge-cloud radius changes association, relaxation, charge-density friction, and form
  factors.                                                                                                                         |
  | **Charge-density mobility**              | Compact high charge density and diffuse charge density must affect mobility differently.
  |
  | **Finite free volume**                   | Hard-sphere volume fraction and void radius must damp mobility as packing increases.
  |
  | **Dense high-salt obstruction**          | High salt / high steric volume must reduce free-carrier translation upstream of the
  readout.                                                                                                                            |
  | **Local obstruction**                    | Obstruction must depend on free volume, ionic strength, additive solvation, carrier size,
  charge density, and solvation support. The implementation already has these factors in `_local_obstruction_factor`.           |
  | **Density / packing scale**              | Measured density or density-derived packing must affect solvent number density and volume
  fraction.                                                                                                                     |
  | **Solvation support**                    | Donor number, acceptor number, and polarizability can reduce or increase local
  obstruction.
  |
  | **Additive solvation support**           | Additive donor/acceptor/HBA/polarizability/volume must affect obstruction and Li
  coordination, not just viscosity.
  |
  | **Anion composition disorder**           | Mixed-anion entropy/disorder must change cation/anion mobility when the mixture is
  compositionally heterogeneous.                                                                                                       |
  | **Solvent composition**                  | EC:DMC:EMC ratio must change viscosity, dielectric, density, donor/acceptor environment,
  packing, and coordination.                                                                                                     |

  The molecular descriptor object includes the raw features needed for these tests: hard-sphere/hydrodynamic/cavity/charge-cloud radii,
  molecular volume, polarizability, donor and acceptor numbers, H-bond counts, dielectric, viscosity, density, Born radius, coordination
  sites, coordination affinity, and ligand-field asymmetry.

  ---

  ## D. Solvent/additive-specific but name-blind effects

  | Effect tests must identify                 | What the test should verify
  |
  | ------------------------------------------ |
  ---------------------------------------------------------------------------------------------------------------------------------------
  ----------------------------------------------------------------------------------------------------------------------------------- |
  | **FEC/VC as coordinating neutral ligands** | The additive must enter as ligand-shell occupancy and Li binding when descriptors imply
  coordination.
  |
  | **FEC viscosity penalty**                  | FEC can lower conductivity through viscosity/free-volume/local obstruction.
  |
  | **FEC solvation support**                  | FEC can also reduce obstruction or alter Li shell structure through donor/acceptor/HBA/
  polarizability descriptors.
  |
  | **FEC–anion competition**                  | FEC must compete with Li–anion coordination and shift SSIP/CIP/free populations.
  |
  | **FEC + bulky-anion anticorrelation**      | Moderate FEC with bulky/delocalized anions must expose a Li–anion cross-current
  coordinate; the test should not require a monotonic FEC penalty.
  |
  | **Additive-separated SSIP**                | A Li–FEC–anion solvent-separated basin must be representable separately from ordinary
  SSIP.
  |
  | **Additive loading basis**                 | Weight fraction vs volume fraction vs solvent fraction conventions must be canonicalized
  and audited.
  |
  | **Additive microviscosity**                | Salt-additive viscosity interaction must affect η, mobility, and rates upstream of σ.
  |
  | **No species-name physics**                | FEC/VC/TFSI/FSI/PF6 names may be registry lookup labels only. Generator states and
  transition logic must use feature IDs. The plan explicitly forbids species-name generator keys, species-name motif labels, molecule-
  specific branches, and salt-name shape hardcoding.  |

  For your motivating example, the test should specifically check this state-local covariance:

  ```text
  D_Li,A | Li shell contains neutral ligand,
           anion has bulky/asymmetric/delocalized feature,
           pair basin = SSIP / additive-separated SSIP / CIP
  ```

  The conductivity primitive should contain:

  [
  D_Q(\alpha)
  ===========

  D_{\rm Li|\alpha}
  +
  D_{\rm A|\alpha}
  ----------------

  2D_{\rm Li,A|\alpha}.
  ]

  If (D_{\rm Li,A|\alpha}<0), the cross term increases σ. A test should fail if the model can only express:

  ```text
  FEC loading -> viscosity increase
  bulky anion -> shape drag
  ```

  but cannot express:

  ```text
  FEC + bulky anion -> changed Li-anion current covariance
  ```

  ---

  ## E. Ion-atmosphere and electrostatic memory effects

  | Effect tests must identify                    | What the test should verify
  |
  | --------------------------------------------- |
  ---------------------------------------------------------------------------------------------------------------------------------------
  ------------------------------------------------------------------------------------------------------------ |
  | **Debye screening length**                    | Ionic strength and dielectric change (\kappa^{-1}), atmosphere resistance, and
  relaxation time.
  |
  | **Electrophoretic drag**                      | Electrophoretic resistance lowers mobility and is separately attributable.
  |
  | **Relaxation drag**                           | Ionic-atmosphere relaxation resistance lowers mobility and is separately
  attributable.
  |
  | **Electrophoretic vs relaxation ownership**   | Audit must split no-atmosphere, electrophoretic-only, relaxation-only, and full-
  atmosphere diffusivities. The plan requires fields like `D_none_alpha`, `D_ep_alpha`, `D_rel_alpha`, `D_full_alpha`, traces, off-
  diagonal norms, and drag factors.  |
  | **Li vs anion diagonal relaxation**           | Audit must split diagonal relaxation into Li-center and anion-center pieces.
  |
  | **Cross relaxation**                          | Li–anion relaxation cross terms must be measurable separately from diagonal terms.
  |
  | **Charge-cloud form factor**                  | Finite charge-cloud radius attenuates relaxation and changes atmosphere resistance.
  |
  | **State geometry form factor**                | SSIP/CIP/aggregate atmosphere projection must depend on internal charge separation,
  not only carrier labels.
  |
  | **Neutral-pair atmosphere cancellation**      | Co-located neutral opposite-charge pairs should cancel long-wavelength atmosphere
  resistance.
  |
  | **Debye-Falkenhagen lifetime gate**           | State lifetime vs atmosphere relaxation time must be a diagnostic; if used, it must
  gate only the declared resistance component.
  |
  | **Atmosphere capture/back-relaxation memory** | Slower exit / stronger capture should increase memory residence and lower effective
  conductivity through the corrector.
  |
  | **Atmosphere bath definition**                | Total-formal vs external-free bath is a diagnostic distinction, not a hidden
  production correction.
  |

  The molecular result object and diagnostics already expose atmosphere-related objects: Debye length, ionic strength, charge-cloud form
  factor by state, friction ratios, bare/electrophoretic/relaxation frictions, countercharge relaxation diffusivity, and atmosphere-
  memory primitive rates/concentrations.

  ---

  ## F. Multi-center transport and internal covariance effects

  | Effect tests must identify               | What the test should verify
  |
  | ---------------------------------------- |
  ---------------------------------------------------------------------------------------------------------------------------------------
  ------------------------------------------------------------- |
  | **State-resolved charged-center tensor** | Conductivity uses (z^T D_\alpha z), not (z_{\rm net}^2D_{\rm COM}).
  |
  | **Neutral Li–anion no-comotion limit**   | A neutral Li–anion state with no co-motion gives (D_Q=D_{\rm Li}+D_{\rm A}).
  |
  | **Perfect comotion neutral pair**        | Equal-diffusivity perfect comotion gives zero charge diffusivity. The plan names this as a
  non-degrading completion gate.                                                                            |
  | **SSIP relative motion**                 | Relative Li–anion motion must be a state-derived covariance/persistence quantity, not a
  formal-charge-only constant.                                                                                 |
  | **CIP comotion**                         | Contact pairs should suppress charge mobility when Li and anion co-move.
  |
  | **Additive-separated relative motion**   | FEC-separated pairs can have different covariance from ordinary SSIP.
  |
  | **Partner switching**                    | Li changing anion partner can carry a charge-displacement moment and memory.
  |
  | **Residence lifetime**                   | Pair, ligand-shell, and aggregate lifetimes must affect (Q), (A,h), or (D_\alpha).
  |
  | **Anion orientation/denticity**          | Bulky/multidentate anion orientation can change mobility, coordination, partner switching,
  and cross-current.                                                                                        |
  | **Bridge/network motifs**                | Multidentate anions must create bridge/network motif concentrations when donor
  functionality and concentration support them. Hostile tests explicitly require failure if such motifs remain absent.  |

  ---

  ## G. Cage, backjump, and structural-memory effects

  | Effect tests must identify                     | What the test should verify
  |
  | ---------------------------------------------- |
  ---------------------------------------------------------------------------------------------------------------------------------------
  --------------------------------------------------------------------------------------------------------------- |
  | **Solvent cage trapping**                      | Caged/mobile substates reduce effective free-carrier transport without terminal
  scaling.
  |
  | **Backjump anticorrelation**                   | Anti-persistent hops increase the Mori corrector and lower σ.
  |
  | **Oriented cage memory**                       | Previous-hop orientation is allowed only when a persistent coordinate survives a
  hop.
  |
  | **Orientation relaxation**                     | Orientation relaxation rate changes memory, not direct σ by itself.
  |
  | **Structural cage exchange**                   | Zero-displacement exchange between mobile/caged states changes memory and occupancy,
  not jump displacement.
  |
  | **Local free-volume stress mode**              | Dense/free-volume modes should be allowed as Mori memory basis functions when they
  couple to current residual.
  |
  | **High-steric compact-anion carbonate defect** | Tests should isolate whether excess σ is owned by free Li translation, free anion
  translation, cage memory, atmosphere, or association. The plan requires event-family attribution and inverse free-translation
  diagnostics for this class of defect.  |

  ---

  ## H. Concentration and high-loading effects

  | Effect tests must identify                     | What the test should verify
  |
  | ---------------------------------------------- |
  ---------------------------------------------------------------------------------------------------------------------------------------
  -------------------------------------- |
  | **Salt molarity**                              | Increasing salt changes carrier density, ionic strength, viscosity, dielectric
  decrement, association, crowding, and atmosphere.                                              |
  | **Flat/overshooting 0.8–1.2 M LiPF6 behavior** | Tests should identify which primitive owns the missing concentration dependence:
  free-carrier mobility, association, atmosphere, cage/backjump, or finite-volume obstruction. |
  | **Finite ionic occupied volume**               | High salt must affect both mobility and finite generator rates. Hostile tests
  require failure if high-salt finite volume is not wired into both.                              |
  | **Ionic-strength crowding**                    | Ionic strength must alter local obstruction and association/crowding stabilization.
  |
  | **Dielectric decrement**                       | Salt loading may alter atmosphere dielectric/screening separately from association
  dielectric.                                                                                |
  | **Jones–Dole / salt viscosity**                | Salt viscosity terms must affect η upstream.
  |
  | **Aggregate onset/saturation**                 | High concentration should shift aggregate populations and free carrier depletion.
  |
  | **Finite-size / packing saturation**           | Packing fraction near max packing should cause loud failure or strong mobility
  suppression, not silent extrapolation.                                                         |

  ---

  ## I. Solvent matrix effects

  | Effect tests must identify         | What the test should verify
  |
  | ---------------------------------- |
  ---------------------------------------------------------------------------------------------------------------------------------------
  -- |
  | **Bulk viscosity**                 | Higher η lowers mobility and rates.
  |
  | **Bulk dielectric**                | Dielectric changes Coulomb association, Born solvation, Debye length, and atmosphere response.
  |
  | **Solvent donor number**           | Stronger donor solvents can stabilize Li solvation and reduce ion pairing.
  |
  | **Solvent acceptor number**        | Acceptors alter solvation support and obstruction.
  |
  | **Solvent polarizability**         | Polarizability contributes to solvation support and local obstruction.
  |
  | **Molecular volume / free volume** | Larger solvent/additive molecules change packing and void radius.
  |
  | **Density**                        | Density changes liquid component number density and packing.
  |
  | **EC/linear carbonate ratio**      | EC:DMC:EMC composition must affect viscosity, dielectric, donor/acceptor environment, and
  packing through descriptors, not solvent names. |
  | **Temperature dependence**         | (T) changes (RT), (k_BT), rates, viscosity if supplied, and conductivity prefactor.
  |

  The current code builds mixture state from dielectric, viscosity, hard-sphere volume fraction, effective solvent radius, mean molecular
  volume, solvent volume fractions, and solvent coordination affinity.

  ---

  ## J. Dataset/audit effects that tests should expose

  | Effect tests must identify               | What the test should verify
  |
  | ---------------------------------------- |
  ------------------------------------------------------------------------------------------------------------------ |
  | **Property label vs trajectory GK gap**  | Separate label error from raw trajectory GK.
  |
  | **Projection gap**                       | Separate raw GK from projected finite-basis GK.
  |
  | **Recipe primitive gap**                 | Separate trajectory-projected primitives from recipe-generated primitives.
  |
  | **D_Q target**                           | Empirical conductivity should be converted to target charge diffusivity for audit.
  |
  | **Uncorrected vs corrected diffusivity** | Report (D_{\rm uncorr}), (D_Q), and (H_{\rm gen}).
  |
  | **Motif population ledger**              | Report free, SSIP, CIP, aggregate, Li-additive, bridge/network populations.
  |
  | **Lifetimes**                            | Report cage, pair, ligand, atmosphere, and aggregate lifetimes.
  |
  | **Top state / top edge attribution**     | Identify which state or event family owns the residual.
  |
  | **Ablation ownership**                   | Association-off, atmosphere-off, viscosity-off, obstruction, cage, and memory ablations
  should be diagnostic only. |
  | **Duplicate/ambiguous empirical rows**   | Inconsistent duplicate formulation labels must be excluded or flagged, not silently
  averaged.                      |
  | **No row dropping**                      | Failed rows must be recorded, not removed from metrics.
  |

  The plan requires empirical audits to decompose predictions into target (D_Q), predicted (D_Q), vehicular, jump/memory, uncorrected,
  memory factor, motif populations, lifetimes, and top contribution rows.

  ---

  ## K. Required descriptor-completeness effects

  A test should fail before conductivity is produced when any needed descriptor is missing. Required descriptor families include:

  ```text
  identity:
    role, charge, molecular weight, density

  size/shape:
    hard-sphere radius, hydrodynamic radius, cavity radius,
    molecular volume, solvent-accessible area, ligand-field asymmetry

  electrostatics:
    charge-cloud radius, dipole, quadrupole, polarizability,
    pure dielectric, Born solvation radius

  coordination/solvation:
    donor number, acceptor number, H-bond donor count,
    H-bond acceptor count, coordination sites,
    coordination affinity

  transport:
    pure viscosity, mixture viscosity, temperature
  ```

  The descriptor backend requires these fields explicitly from user-supplied properties rather than silently filling registry values.

  ---

  ## L. Specific interaction tests I would add first

  These are the highest-value “effect interaction” tests, because they catch exactly the kind of missing physics you are seeing.

  | Interaction                                                     | Test signature
  |                                                                                       |
  | --------------------------------------------------------------- |
  ---------------------------------------------------------------------------------------------------------------------------------------
  ------------------------- | ------------------------------------------------------------------------------------- |
  | **FEC × bulky anion × SSIP**                                    | Increasing neutral ligand loading creates Li-ligand and additive-
  separated SSIP population; state-local (D_{\rm Li,A                                             | \alpha}) changes sign/magnitude; σ is
  allowed to increase despite viscosity increase. |
  | **FEC × PF6 vs FSI/TFSI-like anion**                            | Same additive loading produces different state populations/
  covariances because anion radius, charge cloud, denticity, and ligand-field asymmetry differ.         |
  charge-cloud radius can reduce relaxation drag; tests must identify both channels separately.  |
  |
  | **High LiPF6 concentration × compact anion × carbonate matrix** | If σ remains high/flat, attribution should identify whether free Li
  translation, free anion translation, atmosphere, association, or cage memory owns the error. |
  |
  | **Mixed LiPF6/LiFSI shared-Li**                                 | Free Li pool is shared; adding FSI changes PF6-associated and FSI-
  associated populations without duplicating Li carriers.                                        |
  |
  | **Additive viscosity vs additive coordination**                 | A neutral additive can increase η while also increasing Li
  coordination; tests must show both primitive heads.                                                   |
  |
  | **SSIP relative motion vs CIP comotion**                        | SSIP can increase charge mobility through relative motion; CIP can
  suppress through comotion.                                                                    |
  |
  | **Bridge/network anion functionality**                          | Multidentate anion features create bridge/network populations and
  alter mobility/memory.                                                                         |
  |
  | **Atmosphere diagonal vs cross relaxation**                     | Separate Li diagonal, anion diagonal, and Li–anion cross
  relaxation; do not hide them in one scalar atmosphere factor.                                           |
  |
  | **Solvent-shell identity × Li–anion separation**                | Same Li–anion distance but different shell composition can change
  (D_{\rm Li,A                                                                                   | \alpha}), (Q), and memory.
  |

  ---

  ## Minimal pass/fail checklist

  A conductivity test suite should fail if any of these are true:

  ```text
  FEC changes viscosity but creates no Li-ligand motif.

  Bulky/asymmetric anion changes only scalar radius, not shape,
  charge-cloud, orientation/denticity, or Li-anion covariance.

  SSIP relative/anticorrelated motion is hard-coded from formal charge.

  Neutral multi-charge states use z_net^2 D_COM instead of z^T D z.

  Dual salts count Li once per salt source.

  Pure motif exchange has nonzero charge displacement.

  High salt changes carrier count but not finite volume, viscosity,
  activity, atmosphere, or mobility.

  Ion atmosphere is a scalar terminal correction.

  Poisson/Mori corrector is always zero in systems with cage,
  atmosphere, partner-switch, or state-current memory.

  A generated state label contains FEC, PF6, FSI, TFSI, or other
  species-name physics after descriptor parsing.

  The audit cannot say whether the residual is owned by c_i, Q_ij,
  d/M, D/R, or A/h.
  ```

  That list is the practical definition of “tests can identify the effect.”






